\documentclass[11pt]{article}
\usepackage[margin=1in]{geometry}
\usepackage{amsmath,amssymb,amsthm}
\usepackage{microtype}
\usepackage{enumitem}
\setlist{nosep,leftmargin=*}
\usepackage{booktabs}
\usepackage{xurl}
\usepackage[colorlinks=true,allcolors=blue]{hyperref}
\sloppy
\newtheorem{theorem}{Theorem}
\newtheorem{lemma}[theorem]{Lemma}
\newtheorem{proposition}[theorem]{Proposition}
\newtheorem{corollary}[theorem]{Corollary}
\theoremstyle{definition}
\newtheorem{definition}[theorem]{Definition}
\theoremstyle{remark}
\newtheorem{remark}[theorem]{Remark}
\newcommand{\sharpP}{\#\mathrm{P}}
\newcommand{\FP}{\mathrm{FP}}
\newcommand{\col}{\#\mathrm{col}}
\newcommand{\MGI}{\operatorname{MGI}}
\newcommand{\NAE}{\operatorname{NAE}}
\newcommand{\Q}{\mathbb{Q}}
\newcommand{\Z}{\mathbb{Z}}
\setlength{\parskip}{0.3em}
\title{Counting 2-colourings of 3-regular hypergraphs\\ with Pfaffian incidence graphs}
\author{John Cullen\thanks{Draft of 31 October 2026, prepared with Claude
(Anthropic). The results come from proofs with executable checks in the
public repository \url{https://github.com/jscul/knotlab}; Section~\ref{sec:evidence} lists them, a replay
driver reruns every cited check, and a separately written script recomputes
the headline numbers. No human expert has reviewed this note.}}
\date{}
\begin{document}
\maketitle

\begin{abstract}
Counting the proper 2-colourings of a 3-regular 3-uniform hypergraph is
$\sharpP$-hard (Fan and Cai), and polynomial when its incidence graph is
planar (Valiant). We study the class in between where the incidence graph is
Pfaffian. It contains every hypergraph that cannot be 2-coloured: by Seymour
and by Robertson, Seymour and Thomas, these are exactly the ones whose
incidence matrix is sign-nonsingular. Our results:
(A) for every connected non-2-colourable such hypergraph $H$ and every edge
$L$, the number of 2-colourings of $H-L$ is computable in polynomial time, as
a product over braces of determinants assembled along trisum trees by a
three-term ``trisum law''; (B) with signed incidences, counting stays
polynomial as long as the signs differ from a Pfaffian signing by a twist that
does not branch, because a junction with one untwisted side multiplies the
matchgate identity by $\alpha^2+\alpha\beta+\beta^2$; (C) once twists branch,
counting is $\sharpP$-hard on Pfaffian graphs; (D) it remains $\sharpP$-hard
without any signs; and (E) it remains $\sharpP$-hard when the incidence graph
is a Pfaffian brace. The hardness proofs build crossings inside Pfaffian
graphs: junctions of twisted pieces cross softly, and ladders of plain cube
junctions cross exactly after interpolation. Along the way we find that every
plain planar piece is twisted, and that the Fano plane is the only
non-2-colourable 3-regular 3-uniform hypergraph whose incidence graph is a
brace.
\end{abstract}

\section{Results}\label{sec:results}

A \emph{cubic triple system} $H$ has $n$ points and $n$ lines; every line has
three distinct points and every point lies on three lines. Its incidence graph
$B=B(H)$ is a cubic bipartite graph. A 2-colouring of the points is
\emph{proper} if no line is monochromatic, and $\col(H)$ counts them. In the
language of counting complexity this is $\mathrm{Holant}([0,1,1,0]\mid{=}_3)$,
or not-all-equal 3-SAT with positive literals in which every variable occurs
three times.

Fan and Cai~\cite{FC20} proved $\col$ $\sharpP$-hard; Valiant's holographic
algorithm~\cite{Val08} computes it when $B$ is planar. We consider incidence
graphs that are \emph{Pfaffian}, the class on which perfect matchings can be
counted by a single Pfaffian~\cite{Kas63,Lit75,RST99,McC04}. We also use
\emph{signed} systems: a set $\sigma$ of incidences, at which a line sees the
complement of the point's colour.

\begin{theorem}[counting on Seymour's bound]\label{thm:A}
Let $H$ be a connected cubic triple system that is not 2-colourable, and $L$
one of its lines. Then $\col(H-L)$ can be computed in time polynomial in $n$.
More precisely, $\col(H-L)=2\prod_{K}a(K)$, the product over the braces $K$ of
$B(H)$, where $a(K)=1$ for the two degenerate braces, $a(K)=5$ for the Heawood
graph, $a(K)=|\det|$ of a line--face matrix for a planar brace, and $a(K)$ is
given by the trisum law of Theorem~\ref{thm:law} with two determinants per
planar leaf otherwise. The same holds for every connected cubic signed system
that is not colourable.
\end{theorem}

\begin{theorem}[unbranched twists]\label{thm:B}
Let $T$ be a cubic Pfaffian brace with a trisum tree, $\kappa$ a Pfaffian
signing and $\sigma$ any signing. If no square of the tree has all three of
its branches twisted (Definition~\ref{def:twist}), then $\col(T,\sigma)$ is
computable in polynomial time. The same holds for every connected cubic
bipartite Pfaffian graph through its tight cuts.
\end{theorem}

\begin{theorem}[hardness under signings]\label{thm:C}
Counting $\col(B,\sigma)$ for cubic bipartite Pfaffian graphs $B$ and
signings $\sigma$ is $\sharpP$-hard under polynomial-time Turing reductions.
In the instances used, every junction square has three twisted branches.
\end{theorem}

\begin{theorem}[no signs]\label{thm:D}
Counting $\col(H)$ for cubic triple systems whose incidence graphs are
Pfaffian is $\sharpP$-hard under polynomial-time Turing reductions.
\end{theorem}

\begin{theorem}[braces]\label{thm:E}
Counting $\col(H)$ for cubic triple systems whose incidence graphs are
Pfaffian braces is $\sharpP$-hard under polynomial-time Turing reductions.
\end{theorem}

Theorems~\ref{thm:A} and~\ref{thm:B} describe where the Pfaffian structure
helps: where the signs agree with a Pfaffian signing, counting is a
determinant per planar piece, and twists that run along a path of the trisum
tree do not change that. Theorems~\ref{thm:C}--\ref{thm:E} show that nothing
more survives. Deciding 2-colourability for these systems is polynomial
(Robertson, Seymour and Thomas~\cite{RST99}, McCuaig~\cite{McC04}); counting
is not.

Two structural facts appear on the way (Section~\ref{sec:plain}): every plain
planar piece is twisted, so in a plain trisum brace all twists branch; and
the Heawood graph is the only cubic brace whose plain system cannot be
2-coloured, that is, the Fano plane is the only non-2-colourable cubic triple
system whose incidence graph is a brace. Among plain Pfaffian braces the
polynomial cases we know are therefore planarity and the Fano plane
(Corollary~\ref{cor:picture}).

\paragraph{Provenance.} The source of hardness is Fan and Cai's theorem (we
also give an independent reduction from not-all-equal 3-SAT, which
Creignou and Hermann~\cite{CH96} showed $\sharpP$-complete). Planar
tractability is Valiant's; the structure of Pfaffian braces is due to
Robertson, Seymour and Thomas and to McCuaig; Lemma~\ref{lem:dict} is a
signed form of Seymour's linear-algebra theorem~\cite{Sey74}; Berman's
determinant~\cite{Ber80} is used for planar pieces. Bounded literature
searches found no statement of Theorems~\ref{thm:A}--\ref{thm:E}; folklore
risk is discussed per result in Section~\ref{sec:evidence}.

\section{Preliminaries}\label{sec:prelim}

\paragraph{Signed systems.} A \emph{signed system} is a cubic bipartite graph
$B$ (points and lines) with a set $\sigma\subseteq E(B)$ of negated
incidences. A colouring $c$ of the points is proper if every line $\ell$ sees
the values $c(p)\oplus[p\ell\in\sigma]$ not all equal. Complementing a point's
colour, or all three values a line sees, is a bijection between proper
colourings for $\sigma$ and for $\sigma\triangle\delta(v)$; so $\col(B,\sigma)$
depends only on the class of $\sigma$ modulo cuts.

\paragraph{Pfaffian signings.} A signing $\sigma$ is \emph{Pfaffian} if all
perfect matchings of $B$ contribute with the same sign to the determinant of
the signed biadjacency matrix. A cycle $C$ is \emph{nice} if $B-V(C)$ has a
perfect matching; a signing is Pfaffian if and only if every nice cycle
carries $|C|/2-1$ negations modulo 2, and any two Pfaffian signings of a
connected matching-covered bipartite graph differ by a cut~\cite{Lit75,LP86}.
For planar graphs Kasteleyn's face condition gives one~\cite{Kas63}.

\begin{lemma}[dictionary]\label{lem:dict}
A connected cubic signed system is not colourable if and only if its signing
is Pfaffian. If it is not colourable, deleting any line makes it colourable.
\end{lemma}
For plain systems this is the equivalence of non-2-colourability with
sign-nonsingularity of the incidence matrix (Seymour~\cite{Sey74}; see
\cite[Sect.~1]{RST99}); the signed form is proved in the repository
(\texttt{THREEFOLD-AMPLITUDE-20261030-001C}, Lemma 4). So the non-2-colourable
cubic triple systems are those whose incidence graph is Pfaffian with the
empty signing; we call them \emph{three-fold spaces}.

\paragraph{Braces and trisums.} A \emph{brace} is a connected bipartite graph
in which every matching of two edges extends to a perfect matching; the tight
cut decomposition~\cite{ELP82,LP86} splits a matching-covered bipartite graph
into braces. By Robertson, Seymour and Thomas~\cite{RST99} and
McCuaig~\cite{McC04}, a Pfaffian brace is planar, the Heawood graph, or a
\emph{trisum} of smaller Pfaffian braces: three braces sharing a 4-cycle, the
\emph{square}, glued along it with its four edges deleted. For a cubic
Pfaffian brace $T$ that is neither planar nor the Heawood graph, splitting
repeatedly along squares gives a \emph{trisum tree}: its leaves are planar
braces, its squares are disjoint, each square is a face of the three leaves
containing it, and each square has three \emph{branches}, the components of
the tree minus the square (\texttt{SEYMOUR-COUNT-20261030-001B}, Lemma 1).
A square is written $(c_0,c_1,c_2,c_3)$ with $c_0,c_2$ points and $c_1,c_3$
lines.

\paragraph{The Hadamard form.} With $H=\begin{pmatrix}1&1\\1&-1\end{pmatrix}$
on every edge, a point becomes $2[1,0,1,0]$, a line
$\tfrac14[3,0,-1,0]$ and a negated edge $\mathrm{diag}(1,-1)$, so
\begin{equation}\label{eq:hadamard}
\col(B,\sigma)=\Bigl(\tfrac32\Bigr)^{n}\sum_{\chi}\Bigl(-\tfrac13\Bigr)^{|\chi|/2}(-1)^{|\sigma\cap\chi|},
\end{equation}
the sum over even subgraphs $\chi$ of $B$ (disjoint unions of cycles). Both
constraint functions are matchgate signatures; this is Valiant's planar
algorithm. A 4-ary even function $f$ with legs in the cyclic order of a square
is a matchgate signature if and only if
$\MGI(f)=f_{0000}f_{1111}-f_{1100}f_{0011}+f_{1010}f_{0101}-f_{1001}f_{0110}=0$.

\begin{definition}[twists]\label{def:twist}
Let $\kappa$ be a Pfaffian signing of $B$. The \emph{twist} of a signing
$\sigma$ is $\sigma\triangle\kappa$ modulo cuts. In a trisum tree, the branch
$b$ of a square is \emph{twisted} if $\sigma\triangle\kappa$, restricted to the
edges of the leaves in $b$ (square edges excluded), is not a cut of the graph
they form.
\end{definition}
A branch is untwisted if and only if some sign pattern on its square completes
$\sigma$ to a Pfaffian signing of the branch, and a branch containing a
twisted branch is twisted (\texttt{TWIST-COUNT-20261030-001A}, Lemma 3).

\section{Counting on Seymour's bound}\label{sec:A}

Let $P$ be a \emph{piece}: a connected cubic signed system containing a square,
with a Pfaffian signing in which the square carries an odd number of
negations. Exactly one of $c_1,c_3$ then sees $c_0$ and $c_2$ equal. In a proper
colouring of $P-c_1$, the line $c_3$ sees $c_0$ and $c_2$ different, and its
third point agrees (as seen by $c_3$) with exactly one of them. The
\emph{leaning split} $(\alpha,\beta)$ counts half the colourings of $P-c_1$ in
which it agrees with $c_0$, respectively $c_2$; $\rho=\alpha+\beta$ is the
amplitude $\col(P-L)/2$, the same for every line $L$.

\begin{theorem}[the trisum law]\label{thm:law}
Three pieces in a common gauge, glued along their square, have amplitude
\[
A(T)=\rho_1\rho_2\rho_3-\alpha_1\alpha_2\alpha_3-\beta_1\beta_2\beta_3,
\]
the number of choices of leanings, weighted by the splits, that are not all
equal. Seen from $c_3$ the split is $(\beta,\alpha)$. A piece glued at a
second square $D$ to two pieces $Q,R$ acts on the colourings of its own lines
only through the weights $W_\alpha=\rho_Q\rho_R-\alpha_Q\alpha_R$ and
$W_\beta=\rho_Q\rho_R-\beta_Q\beta_R$.
\end{theorem}
\begin{proof}[Proof outline]
In a common gauge (every piece carries the same odd pattern on the square,
Lemma 1 of \texttt{SEYMOUR-COUNT-20261030-001A}), fix the colours of $c_0,c_2$.
In each piece the square line that sees them equal must be improper inside
the piece, so its third point takes that value; the other square line is then
proper and its third point leans to $c_0$ or $c_2$, with weights $\alpha_i$ or
$\beta_i$. In the trisum, the square lines receive one point from each piece
and must see not-all-equal values; counting gives the law
(\texttt{SEYMOUR-COUNT-20261030-001A}, Theorem 3 and Propositions 4--5).
\end{proof}

\begin{proposition}[a leaf with its children]\label{prop:leaf}
For a planar leaf $\ell$ of a trisum tree, at a square $S$, let $M$ be the
matrix with rows the lines of $\ell$ other than $c_1$ and columns the faces
not at $c_1$, with entry 1 for incidence, except that the rows $d_1,d_3$ of
every other square $D$ of $\ell$ carry $W_\alpha,W_\beta$ on $D$'s two side
faces (in opposite orders) and 1 on $D$'s face. Restricting the row of $c_3$
to its side face through $c_0$, respectively $c_2$, gives matrices
$M_\alpha,M_\beta$, and the split of the branch at $S$ is
$(|\det M_\alpha|,|\det M_\beta|)$.
\end{proposition}
\begin{proof}[Proof outline]
For a planar brace, $|\det|$ of the line--face matrix counts Tutte matchings
with outer triangle at the omitted line, all with one sign (Berman~\cite{Ber80});
weights keep the sign coherence, and Theorem~\ref{thm:law} identifies the
weighted count with the split (\texttt{SEYMOUR-COUNT-20261030-001B},
Proposition 3 and Lemma 4).
\end{proof}

\begin{proof}[Proof of Theorem~\ref{thm:A}]
By Lemma~\ref{lem:dict} and the tight cut decomposition, $\col(H-L)/2$ is a
product of brace amplitudes, independent of $L$ (\texttt{THREEFOLD-AMPLITUDE-20261030-001C},
Corollary 5). The degenerate braces contribute 1 and the Heawood graph 5. A
planar brace contributes the determinant of its line--face matrix. Any other
brace has a trisum tree (Section~\ref{sec:prelim}); its amplitude is the
trisum law at one square, with each branch's split given by
Proposition~\ref{prop:leaf} from the splits one level down. That is two
determinants per leaf, and no signing is needed at any step
(\texttt{SEYMOUR-COUNT-20261030-001B}, Theorem 5).
\end{proof}

\begin{remark}
The checks compare the algorithm with direct counts on 67 braces, on
three-fold spaces of 22 to 118 points (for example amplitudes $125$ for a
22-point space and $3906250000$ for a 90-point one, an amplitude being
$\col(H-L)/2$), and on trisums of omnitruncated polyhedra up to 480 vertices; trisums of three omnitruncated $9\times 9$ grids (1720
vertices) are counted in seconds, where variable elimination would need tables
of width 30.
\end{remark}

\section{Twists}\label{sec:B}

Let $T$ be a cubic Pfaffian brace with a trisum tree and $\sigma$ a signing.
For a square $S$ and a branch $b$, the \emph{branch tensor}
$F_{S,b}(z_0,z_1,z_2,z_3)$ counts the colourings of the branch's own points
that make its own lines proper, where $z_0,z_2$ are the colours of $c_0,c_2$
and $z_1,z_3$ those of the neighbours of $c_1,c_3$ in $b$. Seen from a third
branch, the \emph{junction} of two closed branches $Q,R$ is
\[
W(x_0,z_1,x_2,z_3)=\sum F_Q(x_0,q_1,x_2,q_3)F_R(x_0,r_1,x_2,r_3)\,
\NAE(z_1,q_1,r_1)\,\NAE(z_3,q_3,r_3),
\]
with incidence signs complementing the corresponding arguments. In the
Hadamard basis, branch tensors are even, and a single planar leaf is a
matchgate signature under any signing.

\begin{theorem}[the junction is a similitude]\label{thm:sim}
If $Q$ is untwisted with split $(\alpha,\beta)$, then for every $R$
\[
\MGI(\hat W)=\pm(\alpha^2+\alpha\beta+\beta^2)\,\MGI(\hat F_R),
\]
the sign given by four incidence signs at the junction. In particular $W$ is a
matchgate signature if and only if $F_R$ is.
\end{theorem}
\begin{proof}[Proof outline]
An untwisted branch has the \emph{split form}: a matchgate signature
determined by $(\alpha,\beta)$, an odd pattern on the square and the leg signs
(\texttt{TWIST-COUNT-20261030-001A}, Lemma 3). The identity is then a
polynomial identity in $\alpha,\beta$ and the eight even entries of
$\hat F_R$, verified symbolically for all eight odd patterns and all four
choices of leg signs (ibid., Theorem 4).
\end{proof}

The multiplier is $|\alpha+\beta e^{i\pi/3}|^2$, a norm in the Eisenstein
integers; it is positive, so an untwisted side never moves a junction off the
matchgate class.

\begin{proof}[Proof of Theorem~\ref{thm:B}]
If no square has three twisted branches, some leaf $r$ sees at most one
twisted branch among the two closed branches of every square
(\texttt{TWIST-COUNT-20261030-001A}, Lemma 5). Rooted at $r$, every junction
has an untwisted closed branch, so by Theorem~\ref{thm:sim} and induction
every junction is a matchgate signature, and each leaf is a planar matchgate
computation (FKT). Tight cuts contribute three-edge gadgets, which are always
matchgates (ibid., Theorem 6).
\end{proof}

When twists branch the mechanism stops: junctions of two twisted pages leave
the matchgate class. They still cannot cross wires sharply.

\begin{proposition}[soft crossings]\label{prop:soft}
(i) With nonnegative branch tensors no junction is an exact crossover
$[x_0=x_2\oplus a][z_1=z_3\oplus b]$: the square's lines forbid one pattern
only, so on any slice confined to one parity class of $(z_1,z_3)$ the junction
pins $z_1$ (\texttt{TWIST-COUNT-20261030-001B}, Theorem 2).
(ii) Junctions of two twisted truncated-octahedron pages include 256 that are
\emph{crossover-type} (zero turning entries in the Hadamard basis) with all
four straight entries nonzero, for example $(40000,-1112,3336,-192)$
(ibid., Proposition 3). We call such a junction a \emph{soft crossing}.
\end{proposition}

\section{Hardness under signings}\label{sec:C}

A \emph{dead cap} is the cube minus an edge, attached to an outside line and
an outside point: its count is 3 whatever they see, so a line with one capped
leg acts on its other two legs by $3M$, $M=\begin{pmatrix}1&2\\2&1\end{pmatrix}$.
A \emph{$k$-snake} is $k$ soft crossings in series, the wires joined between
junctions by capped lines and points.

\begin{lemma}[snakes and the lattice condition]\label{lem:snake}
The $k$-snake built from the soft crossing $\sigma_1=(40000,-1112,3336,-192)$
is crossover-type with straight entries $\sigma_1(s)\Lambda(s)^{k-1}$, where
\[
\Lambda=(810000,\ 7506,\ -22518,\ -432)=(2^43^45^4,\ 2\cdot3^3\cdot139,\ -2\cdot3^4\cdot139,\ -2^43^3).
\]
The valuations at $5$, $139$ and $3$ determine any count vector $n\in\mathbb N^4$
of given total from $\prod_s\Lambda(s)^{n_s}$, so the products
$\prod_s\Lambda(s)^{2n_s}$ are distinct.
\end{lemma}
\begin{proof}
The multiplier per step is $\tfrac94\mu(y_A)\mu(y_B)\sigma_1(s)$ with
$\mu=(3,-1)$ the eigenvalues of $M$; the valuations are $v_5=(4,0,0,0)$,
$v_{139}=(0,1,1,0)$ and $v_3=(4,3,4,3)$, which give $n_{00}$, then $n_{01}+n_{10}$
and $n_{10}$ (\texttt{SOFT-CROSS-20261031-001A}, Lemmas 2--3).
\end{proof}

\begin{theorem}[universality inside Pfaffian graphs]\label{thm:univ}
Let $G$ be a cubic signed system drawn in the plane with $m$ crossings. Replace
every edge by a capped path with $k_E$ pairs and every crossing by a
$k_S$-snake. The resulting instances are Pfaffian, and $\col(G)$ is computed
in polynomial time from their counts for $k_E$ in an interval of length
$|E(G)|+1$ and $k_S=1,3,\dots,2N-1$, $N=\binom{m+3}{3}$.
\end{theorem}
\begin{proof}[Proof outline]
The host (everything outside the pages) is planar and every junction square
is a face of it and of its pages, so Kasteleyn signings of all pieces,
switched to a common odd pattern on each square, glue to a non-colourable,
hence Pfaffian, signing (Theorem~\ref{thm:law}'s gauge lemma and
Lemma~\ref{lem:dict}). Expanding long edges into $9^kP_+ +(-3)^kP_-$ and snakes
into their straight states gives a sum over two families of distinct nodes
(Lemma~\ref{lem:snake}); a two-parameter Vandermonde system determines every
coefficient, and replacing each snake by the compensated crossover
$[z_1=z_3]\,w(z_0,z_2)$ with $w=\begin{pmatrix}-1&2\\2&-1\end{pmatrix}$ gives
$9^m\col(G)$ (\texttt{SOFT-CROSS-20261031-001A}, Theorem 5).
\end{proof}

\begin{proof}[Proof of Theorem~\ref{thm:C}]
Counting the colourings of plain cubic triple systems is $\sharpP$-hard: this
is Fan and Cai's theorem for $\mathrm{Holant}([0,1,1,0]\mid{=}_3)$, and the
repository also derives it from positive not-all-equal 3-SAT~\cite{CH96} by
joining the copies of a variable with capped paths, which act by
$9^k(P_++(-\tfrac13)^kP_-)$ and interpolate to equalities
(\texttt{SOFT-CROSS-20261031-001B}, Theorem 3). Theorem~\ref{thm:univ} reduces
these to Pfaffian instances. Both pages of every junction are twisted, a
cut restricts to a cut, and the host side of a junction in a snake contains
other junctions' pages, so all three branches are twisted
(ibid., Theorem 4).
\end{proof}

\section{No signs}\label{sec:plain}

\begin{lemma}[signs are path lengths]\label{lem:long}
Replace each of the $L$ negated incidences $(p,\ell)$ of a signed system by
an unsigned capped path with $k$ pairs, acting by $9^k(P_++tP_-)$ with
$t=(-\tfrac13)^k$. The count is $9^{kL}q(t)$ for a polynomial $q$ of degree at
most $L$, and $q(-1)$ is the signed count while $q(1)$ is the count with the
signs dropped.
\end{lemma}
\begin{proof}
At $t=-1$ the path acts by $P_+-P_-=\begin{pmatrix}0&1\\1&0\end{pmatrix}$, the
negation; the nodes $(-\tfrac13)^k$, $k=1,\dots,L+1$, are distinct
(\texttt{PLAIN-COUNT-20261031-001A}, Lemma 1).
\end{proof}

\begin{proof}[Proof of Theorem~\ref{thm:D}]
Capped paths drawn inside the planar pieces keep them planar, so the
certificate of Theorem~\ref{thm:univ} survives. Switching inside the pages
leaves six negations per junction. Lemma~\ref{lem:long} then removes all
signs from the instances of Theorem~\ref{thm:C}
(\texttt{PLAIN-COUNT-20261031-001A}, Theorem 4).
\end{proof}

\begin{proposition}[plain pages are twisted]\label{prop:pages}
(i) Every 2-connected planar cubic bipartite graph has at least six square
faces, and at most five faces meet a given square face. (ii) Hence every
\emph{plain page} (such a graph with a junction square face, no signs) is
twisted. (iii) In a plain cubic Pfaffian brace that is a trisum, every square
has three twisted branches, so Theorem~\ref{thm:B} applies to none.
\end{proposition}
\begin{proof}
(i) $\sum_f(6-|f|)=12$ by Euler's formula, and only squares contribute
positively; a face meeting a square $S$ at a vertex contains an edge of $S$.
(ii) A square face $Q$ disjoint from $S$ is a cycle of the page's own edges, a
Kasteleyn signing has $|Q|/2-1=1$ negation on it, and the plain signing has
none, so they differ by a non-cut. (iii) Every branch contains a leaf of the
tree, a plain page, and twistedness is inherited
(\texttt{PLAIN-COUNT-20261031-001B}, Theorem 2 and Corollary 3).
\end{proof}

\begin{proposition}[the Fano plane among braces]\label{prop:fano}
In a plain non-colourable system every nice cycle has length $2$ modulo $4$.
Consequently the Heawood graph is the only cubic brace whose plain system is
not 2-colourable: the Fano plane is the only non-2-colourable cubic triple
system whose incidence graph is a brace.
\end{proposition}
\begin{proof}
The empty signing is Pfaffian (Lemma~\ref{lem:dict}), so a nice cycle $C$
needs $|C|/2-1\equiv0$. In a brace every 4-cycle is nice, so a plain
non-colourable brace has no 4-cycle. Planar cubic bipartite graphs have
square faces (Proposition~\ref{prop:pages}), and a trisum contains a square
face of a leaf away from its junction; by Robertson, Seymour and Thomas only
the Heawood graph remains (\texttt{PLAIN-COUNT-20261031-001B}, Corollary 5).
\end{proof}
This is probably implied by the description of minimal non-2-colourable
square hypergraphs as splices of Fano planes~\cite{RST99}; the argument above
is short and direct.

\section{Braces}\label{sec:E}

The instances so far contain 2-edge cuts (caps and capped paths), so they are
not braces, and a single junction of plain pages is never a crossing.

\paragraph{Ladder sites.} Replace a crossing of edges $e$ and $f$ by $k$
junction squares, each glued by trisum to two plain cubes. Wire $A$ runs from
$e$'s point end through $l_s$, then $c_0\to c_2$ of each junction; wire $B$
runs from $f$'s line end through $\pi_s$, then $c_1\to c_3$ of each junction.
The start rung $l_s\pi_s$ and a rung $\lambda_j\pi_j$ between consecutive
junctions close every leg; there are no caps.

\begin{lemma}[the site is a matrix power]\label{lem:power}
As a map from (colour of $e$'s point end, colour of $\pi_s$) to (colour seen by
$e$'s line end, colour of $f$'s point end) the site acts by $(NW)^k$, where $W$
is the junction and $N[(z_2,z_3),(a,b)]=\NAE(z_2,a,b)[z_3=b]$ the rung.
Restoring $e$ and $f$ is the identity, the term $k=0$.
\end{lemma}

\begin{lemma}[the cube junction]\label{lem:cube}
For two plain cubes,
\[
NW=\begin{pmatrix}2&3&2&5\\8&4&8&9\\9&8&4&8\\5&2&3&2\end{pmatrix},\qquad
\det(xI-NW)=(x^2-19x-1)(x^2+7x+13).
\]
The roots are $\varepsilon,-1/\varepsilon$ with $\varepsilon=(19+\sqrt{365})/2$ a
unit of norm $-1$, and $\pi,\bar\pi$ with $\pi=(-7+i\sqrt3)/2$ an Eisenstein
prime of norm 13. For count vectors $n\ne n'$ of equal total,
$\prod\lambda_i^{2n_i}\ne\prod\lambda_i^{2n'_i}$.
\end{lemma}
\begin{proof}
If two such products agree, $\varepsilon^{2(a-a')}$ (in $\Q(\sqrt{365})$)
equals $\pi^{2(b'-b)}\bar\pi^{2(c'-c)}$ (in $\Q(\sqrt{-3})$), so both are
rational; $\pi/\bar\pi$ is not a root of unity, so $b'-b=c'-c$, and a rational
power of the unit $\varepsilon$ is trivial
(\texttt{PLAIN-BRACE-20261031-001A}, Lemmas 2--3).
\end{proof}

Consequently, with $Z(k)$ the count when a single site has $k$ junctions, the
count with the crossing restored is
$\sum_{j=0}^{3}f_jZ(2j+1)$, where
$f=(-\tfrac{53451621}{2804438},\tfrac{521319}{200317},-\tfrac{23524}{200317},\tfrac{853}{2804438})$
is the unique cubic with $xf(x^2)\equiv1$ modulo the characteristic
polynomial; for $m$ sites sharing an odd $k$, $\binom{m+3}{3}$ counts suffice.
Only odd $k$ gives a planar host. For $K_{3,3}$ drawn with one crossing, the
plain instances with $k=1,3,5,7$ have 20, 48, 76 and 104 vertices and counts
80, 31500, 11459396 and 4159886188, and the formula returns $\col(K_{3,3})=6$.

\begin{lemma}\label{lem:braces}
(i) A connected cubic bipartite graph on at least six vertices that is
cyclically 4-edge-connected is a brace. (ii) A site with $k$ junctions is the
ladder $P_{3k}\times K_2$ with its ports at the corners. (iii) Let a connected
cubic bipartite graph $G$ be drawn with its vertices in distinct faces of a
cyclically 4-edge-connected planar cubic bipartite frame $F$, all edges in
general position. Then every bond with at most three edges of the planarized
drawing isolates one vertex, and the host obtained by replacing every crossing
by a site is cyclically 4-edge-connected. (iv) Gluing two cubes onto a
4-cycle of a cyclically 4-edge-connected cubic graph, the cycle's edges
deleted, keeps it cyclically 4-edge-connected.
\end{lemma}
\begin{proof}[Proof outline]
(i) If $|N(X)|\le|X|+1$ for some $X$ on one side with $2\le|X|\le n-2$, the cut
around $X\cup N(X)$ has three edges and both sides contain cycles; the Hall-type
characterization of braces~\cite{LP86} concludes. (ii) The rungs are the start
rung, the two square edges of each junction and the rungs between junctions.
(iii) A closed curve crossing at most three edges has at most one frame vertex
inside (the frame is cyclically 4-edge-connected), and the one-vertex-per-face
condition leaves room for at most one vertex of $G$; sites are ladders, so a
site cut through costs two edges and its pieces with a cycle have four boundary
edges. (iv) The twelve vertices of the junction are 2-edge-connected, their
subsets with a cycle have at least four boundary edges, and the host minus the
square has no bridge (\texttt{PLAIN-BRACE-20261031-001B}, Lemmas 1--5).
\end{proof}

\begin{proof}[Proof of Theorem~\ref{thm:E}]
Take a hard plain source $G$ (Theorem~\ref{thm:C}'s proof), connected, and draw
it with its vertices in distinct square faces of a prism $C_{2n}\times K_2$,
$2n\ge|V(G)|$, drawn by Tutte's method~\cite{Tut63}. Replace every crossing by
a site with $k=2t+1$ junctions. By Lemma~\ref{lem:braces} each instance is
cyclically 4-edge-connected, hence a brace; it is Pfaffian by the certificate
of Theorem~\ref{thm:univ} (every 4-cycle of the host is a face); it carries no
sign. By Lemma~\ref{lem:cube} the counts for $t<\binom{m+3}{3}$ determine
$\col(G\sqcup F)=\col(G)\col(F)$, and $\col(F)\ne0$ because $F$ has nice
4-cycles (Proposition~\ref{prop:fano})
(\texttt{PLAIN-BRACE-20261031-001B}, Theorem 7).
\end{proof}

\begin{corollary}\label{cor:picture}
Among cubic triple systems whose incidence graphs are Pfaffian braces, the
planar ones count in polynomial time (Valiant), the Fano plane is a single
system, and the class is $\sharpP$-hard. The Fano plane is the only one
without twists.
\end{corollary}

\section{Remarks and open questions}\label{sec:remarks}

\paragraph{The picture.} On Pfaffian incidence graphs the count is governed by
the sign class: the Pfaffian class gives one determinant per planar piece
(Theorem~\ref{thm:A}), twists along a path of the trisum tree change nothing
(Theorem~\ref{thm:B}), and branching twists make the count universal
(Theorems~\ref{thm:C}--\ref{thm:E}). Two twisted sides at a square pass as a
similitude when the third side is untwisted; three twisted sides make a soft
crossing; a ladder of plain squares crosses exactly.

\paragraph{An Ising reading.} By~\eqref{eq:hadamard}, $\col(B,\sigma)$ is an
Ising partition function on $B$ at the imaginary coupling with $x^2=-\tfrac13$,
the signing giving the signs of the couplings; Pfaffian signings are its
free-fermion sign patterns. We have not used this reading.

\paragraph{Open questions.}
\begin{itemize}
\item \emph{Few branched squares.} An even 4-ary signature with a nonzero
entry at $0000$ is a matchgate signature plus a multiple of the indicator of
$1111$. This suggests an algorithm exponential only in the number of squares
with three twisted branches; we have not written it down.
\item \emph{Approximation.} Approximate counting on these classes is untouched;
the degree-3 and edge-size-3 setting lies between the known algorithmic and
hardness thresholds for hypergraph colourings.
\item \emph{Deficiency one.} Theorem~\ref{thm:A} uses that $H$ sits on
Seymour's bound ($n$ lines on $n$ points). Systems with one line fewer are
open.
\end{itemize}

\section{Evidence and replay}\label{sec:evidence}

Every theorem is proved in the repository, with a \texttt{check.py} that
recomputes the cited facts and controls that must fail. Proof directories are
\texttt{proofs/<ticket>/<ticket>-P001\_<slug>/}.

\begin{center}
\footnotesize
\begin{tabular}{@{}p{3.3cm}p{6.2cm}p{5.2cm}@{}}
\toprule
Result & Proved in & Novelty, folklore risk\\
\midrule
Lemma~\ref{lem:dict} & \texttt{THREEFOLD-AMPLITUDE-20261030-001C} & signed form of Seymour's theorem; high\\
Theorem~\ref{thm:law}, Prop.~\ref{prop:leaf} & \texttt{SEYMOUR-COUNT-20261030-001A}, \texttt{SEYMOUR-COUNT-20261030-001B} & no statement found; moderate\\
Theorem~\ref{thm:A} & \texttt{SEYMOUR-COUNT-20261030-001B} & no statement found; moderate\\
Theorem~\ref{thm:sim}, Theorem~\ref{thm:B} & \texttt{TWIST-COUNT-20261030-001A} & no statement found; low--moderate\\
Proposition~\ref{prop:soft} & \texttt{TWIST-COUNT-20261030-001B} & exact computation; low\\
Lemma~\ref{lem:snake}, Theorem~\ref{thm:univ} & \texttt{SOFT-CROSS-20261031-001A} & standard method, new setting; moderate\\
Theorem~\ref{thm:C} & \texttt{SOFT-CROSS-20261031-001B}, \texttt{SOFT-CROSS-20261031-001C} & source known (Fan--Cai); moderate\\
Lemma~\ref{lem:long}, Theorem~\ref{thm:D} & \texttt{PLAIN-COUNT-20261031-001A} & standard method, new setting; moderate\\
Props.~\ref{prop:pages}, \ref{prop:fano} & \texttt{PLAIN-COUNT-20261031-001B} & Prop.~\ref{prop:fano} likely known; high\\
Lemmas~\ref{lem:power}, \ref{lem:cube} & \texttt{PLAIN-BRACE-20261031-001A} & no statement found; low--moderate\\
Lemma~\ref{lem:braces}, Theorem~\ref{thm:E} & \texttt{PLAIN-BRACE-20261031-001B} & no statement found; moderate\\
\bottomrule
\end{tabular}
\end{center}

\paragraph{Computations.} Among the checks: the similitude identity
symbolically for all 32 pattern and sign choices; 13120 junctions of cube and
6-prism pages, none a crossover; 256 soft crossings among truncated-octahedron
pages; $K_{3,3}$ recovered from Pfaffian instances under three signings
($54,18,18=9\cdot(6,2,2)$); six formulas recovered through the source
reduction on graphs up to 2710 vertices; two formulas recovered end to end
through 44 certified Pfaffian instances; $K_{3,3}$ recovered from 100 counts of
plain Pfaffian instances, the largest with 18504 vertices; 454 plain pages
over 33 planar graphs, all twisted; woven brace instances of 214 to 1078
vertices, certified Pfaffian, with a site recovered exactly inside them.

\paragraph{Replay.} \texttt{replay.py} reruns every cited check in a
temporary copy and compares each fresh result with the frozen one, ignoring
timings. Two checks depend on Python's per-run hash order: one labels the two
weights of a child square in that order, the other draws random signings.
There the replay accepts equal polynomials up to exchanging those weights, and
equal numbers of signings of each kind with no failures; every other field
must match exactly; \texttt{independent.py}, written separately without the repository's
libraries, recomputes the soft crossing, the cube junction and its
characteristic polynomial, the $K_{3,3}$ counts through plain braces, sign
interpolation, the similitude identity, and Theorem~\ref{thm:A} on
compositions of Fano planes; \texttt{check.py} ties every number quoted here to
those runs. The review package lists the claims, the replay steps and
questions for a reviewer.

\begin{thebibliography}{99}
\bibitem{Ber80} K.~A.~Berman, \emph{A proof of Tutte's trinity theorem and a
new determinant formula}, SIAM J. Algebraic Discrete Methods 1 (1980) 64--69.
\bibitem{CH96} N.~Creignou and M.~Hermann, \emph{Complexity of generalized
satisfiability counting problems}, Inform. and Comput. 125 (1996) 1--12.
\bibitem{CLX10} J.-Y.~Cai, P.~Lu and M.~Xia, \emph{Holographic algorithms with
matchgates capture precisely tractable planar \#CSP}, FOCS 2010.
\bibitem{ELP82} J.~Edmonds, L.~Lov\'asz and W.~R.~Pulleyblank, \emph{Brick
decompositions and the matching rank of graphs}, Combinatorica 2 (1982)
247--274.
\bibitem{FC20} A.~Z.~Fan and J.-Y.~Cai, \emph{Dichotomy result on 3-regular
bipartite non-negative functions}, arXiv:2011.09110 (2020).
\bibitem{Kas63} P.~W.~Kasteleyn, \emph{Dimer statistics and phase
transitions}, J. Math. Phys. 4 (1963) 287--293.
\bibitem{Lit75} C.~H.~C.~Little, \emph{A characterization of convertible
$(0,1)$-matrices}, J. Combin. Theory Ser. B 18 (1975) 187--208.
\bibitem{LP86} L.~Lov\'asz and M.~D.~Plummer, \emph{Matching Theory},
North-Holland, 1986.
\bibitem{McC04} W.~McCuaig, \emph{P\'olya's permanent problem}, Electron. J.
Combin. 11 (2004) R79.
\bibitem{RST99} N.~Robertson, P.~D.~Seymour and R.~Thomas, \emph{Permanents,
Pfaffian orientations, and even directed circuits}, Ann. of Math. 150 (1999)
929--975.
\bibitem{Sey74} P.~D.~Seymour, \emph{On the two-colouring of hypergraphs},
Quart. J. Math. Oxford 25 (1974) 303--312.
\bibitem{Tut63} W.~T.~Tutte, \emph{How to draw a graph}, Proc. London Math.
Soc. 13 (1963) 743--767.
\bibitem{Val08} L.~G.~Valiant, \emph{Holographic algorithms}, SIAM J. Comput.
37 (2008) 1565--1594.
\bibitem{VY89} V.~V.~Vazirani and M.~Yannakakis, \emph{Pfaffian orientations,
0-1 permanents, and even cycles in directed graphs}, Discrete Appl. Math. 25
(1989) 179--190.
\end{thebibliography}
\end{document}
